\section{Evaluation}
\label{sec:evaluation}

Each research question is answered by the live experiments that bear on it
most directly; the remaining attack scenarios, the prover-backend
comparison, and the Bitcoin-reorg boundary test appear in
Appendices~\ref{app:attacks}, \ref{app:backend}, \ref{app:aggregation} and~\ref{app:reorg}. Unless stated otherwise, all
numbers are live measurements on Testbed~A, the seven-validator testbed of
Section~\ref{sec:setup}; its configuration and the FSM parameters are in
Appendix~\ref{app:setup} (Tables~\ref{tab:environment} and~\ref{tab:fsm-params}), and
metrics follow Section~\ref{subsec:metrics}. Table~\ref{tab:provenance} in Appendix~\ref{app:setup} states for each result whether it is live, in-process, model-based or offline, what one replicate is, and which recovery mode the runs used.

\subsection{RQ1: Safety Under Byzantine Behavior}
\label{subsec:rq1}

By safety we mean agreement among honest validators (identical FSM state and
AppHash, aggregate within the honest range, forged or stale receipts rejected)
and the lock invariant on committed blocks: no withdrawal-marked transaction
commits while the state is \textsc{Sovereign} or \textsc{Recovering}. We do not
claim that cross-chain assets are protected; the lock acts only once the FSM has
reached a locked state (Section~\ref{subsec:rq2}), and the prototype has no bridge
(Section~\ref{sec:discussion}).

\textbf{Single Byzantine proposer.} The most direct attack on the extended
consensus object is a proposer that claims an FSM state not implied by the
certificate. One of the seven validators ($n{=}7$, 14.3\% of the voting
power) was configured to do so. Every honest validator
recomputed $\delta(s^{h-1}, \mathsf{Agg}(C^{h}))$, rejected the forged claim
in \texttt{ProcessProposal}, and the chain committed only honest proposals,
with zero AppHash divergence. The same outcome held for the nine further
single-actor attacks of Appendix~\ref{app:attacks} (eclipse, Sybil, data
withholding, forged Bitcoin receipt, withdrawal in \textsc{Sovereign},
censorship, a combined attack, double-signing, and timeout flooding).

The analysis of colluding Byzantine validators ($f{=}2$) is in Appendix~\ref{app:extended}; the honest validators agreed on the AppHash at every height in all cases.

The function-level checks of the aggregation are in Appendix~\ref{app:extended}.

\subsection{RQ2: Liveness Under Peripheral Degradation}
\label{subsec:rq2}

\textbf{Engram versus the strict receipt-gated baselines.} The two receipt-gated
baselines (Section~\ref{subsec:baselines}) enforce Bitcoin and DA receipts at
every height, one reading them through the vote extensions (in-band) and one
directly at each validator (out-of-band observation); the soft baseline keeps
Engram's FSM but never locks withdrawals. We ran each system through three faults: Bitcoin checkpoint submission suspended,
all seven Celestia bridges stopped, and both. A trial is 60\,s healthy, 180\,s
of fault and 120\,s of recovery, with a withdrawal-marked transaction sent every
2\,s; each cell of Table~\ref{tab:rq2} has three trials. Both
receipt-gated baselines are redeployed for every trial, because the in-band one
does not recover on its own in our implementation (below).
Engram and the soft baseline recover through proof-only recovery with an
unproven-tail bound of 128 blocks (Section~\ref{subsec:reanchoring}); the
receipt-gated baselines have no recovery mode, since their FSM stays \textsc{Anchored}.

\begin{table}[tbp]
\caption{Liveness and withdrawals under peripheral degradation (180\,s fault,
three runs per cell, mean and range). $n$ is the number of trials; Blocks is the count committed in the fault window, Max round the highest round at which a proposal was decided and Longest gap the longest interval between consecutive blocks in the window. Rate is the fault-window block rate as a
share of the same trial's healthy rate; ``No resume'' counts trials in which the
chain committed no block in the 120\,s after the fault ended. A gap entered as
$>$$x$ is a chain that never resumed: $x$ is the time from its last block to the end
of the trial, a lower bound of a gap that never closed. The in-band gate is the receipt-gated baseline whose receipts come
through the vote extensions, the out-of-band gate the
receipt-gated one in which each validator reads the layers itself;
Soft is Engram without the withdrawal lock. The last two columns count withdrawals committed
after and before the moment Engram's circuit breaker locks, 125\,s (checkpoint
pause), 40\,s (DA outage) and 40\,s (both) after the fault begins; the same
window is applied to every system. Engram and the soft baseline recover through
proof-only recovery with an unproven-tail bound of 128 blocks.}
\label{tab:rq2}
\centering
\footnotesize
\setlength{\tabcolsep}{3pt}
\resizebox{\linewidth}{!}{% Generated by scripts/e11_hard_gated_baseline/make_latex_table.py -- do not edit.
\begin{tabular}{@{}llrrrrrrrr@{}}
\toprule
 & & & \multicolumn{2}{c}{Liveness in the fault} & & & & \multicolumn{2}{c}{Withdrawals committed} \\
\cmidrule(lr){4-5}\cmidrule(lr){9-10}
\textbf{Fault} & \textbf{System} & $n$ & Blocks & Rate & Longest gap (s) & Max round & No resume & After lock & Before lock \\
\midrule
\multirow{4}{*}{Checkpoint pause} & In-band gate & 3 & 81 (79--82) & 64 (56--78)\% & $>$181 & 13 & 3/3 & 0 & 59 \\
 & Out-of-band gate & 3 & 89 & 61 (60--63)\% & 6.7 (6.6--6.8) & 12 & 0/3 & 3 & 59 (58--59) \\
 & Engram & 3 & 142 (139--145) & 96 (94--98)\% & 6.3 (6.2--6.4) & 2 (2--3) & 0/3 & 0 & 58 (55--59) \\
 & Soft & 3 & 143 (142--144) & 97 (96--98)\% & 6.4 (6.3--6.4) & 2 (2--3) & 0/3 & 28 (27--29) & 60 (58--61) \\
\midrule
\multirow{4}{*}{DA outage} & In-band gate & 3 & 26 (23--28) & 24 (22--25)\% & $>$264 (261--267) & 21 (21--22) & 3/3 & 5 (3--7) & 18 (17--19) \\
 & Out-of-band gate & 3 & 29 (29--30) & 20 (19--21)\% & 6.7 (6.6--6.8) & 21 (20--21) & 0/3 & 8 (7--8) & 19 (18--19) \\
 & Engram & 3 & 150 (149--151) & 100 (99--103)\% & 1.4 & 0 & 0/3 & 0 & 18 (18--19) \\
 & Soft & 3 & 149 (146--152) & 98 (96--100)\% & 1.4 (1.3--1.4) & 0 & 0/3 & 72 (72--73) & 19 \\
\midrule
\multirow{4}{*}{Both} & In-band gate & 3 & 27 (24--28) & 17 (16--18)\% & $>$262 (259--266) & 21 & 3/3 & 6 (3--7) & 19 (18--19) \\
 & Out-of-band gate & 3 & 30 (27--33) & 20 (17--22)\% & 6.7 (6.7--6.8) & 20 (20--21) & 0/3 & 9 (7--10) & 18 (18--19) \\
 & Engram & 3 & 149 (148--151) & 97 (95--98)\% & 1.4 (1.3--1.4) & 0 & 0/3 & 0 & 17 (15--18) \\
 & Soft & 3 & 150 (149--151) & 97 (96--97)\% & 1.3 & 0 & 0/3 & 72 (72--73) & 19 \\
\bottomrule
\end{tabular}
% lock times used for the two withdrawal columns: checkpoint pause 125\,s, da outage 40\,s, both 40\,s
}
\end{table}

Two observations should be kept apart. \emph{During the fault}, the receipt-gated
baselines commit a fraction of their blocks or stop, which the validity predicate
predicts and the experiment measures. The out-of-band baseline commits 20\%
(29--30 blocks) of its healthy rate under a DA outage, 20\% (27--33 blocks)
under the combined fault and 61\% (89 blocks) under checkpoint suspension, with a
longest gap of 6.6--6.8\,s and rounds up to 21 (12 under checkpoint suspension),
against \resRQTwoEngram{} for Engram. The in-band baseline does not slow but halts: it commits
79--82 blocks under checkpoint suspension, 23--28 under a DA outage and 24--28 under the combined
fault, all before it stops, 120--121\,s after a checkpoint fault begins and 44--52\,s after a DA or
combined fault begins, and it commits no block again until the end of the trial (the open gap is at least
181\,s and 259\,s). The fault-window rate in the table, 56--78\%, 22--25\% and 16--18\% of the healthy rate,
is therefore the count of blocks before the halt divided by 180\,s; proposals reach round 13 under
checkpoint suspension and 21--22 otherwise. Its healthy
rate differs between trials (0.55 to 0.82 blocks/s), which is why that rate
range is wider than the block counts. In the four trials with the lowest healthy
rate (0.55 to 0.60) the 60\,s healthy window contains three or four blocks decided after
round 0, each preceded by a gap of about 6.5\,s; the other five in-band trials and all Engram trials contain none. The Bitcoin height gap of
a healthy chain reaches the receipt tolerance once per checkpoint cycle
(Appendix~\ref{app:envelope}), which is consistent with this, but we did not
record the rejected proposals (Section~\ref{subsec:limitations}).
Because the rate of Table~\ref{tab:rq2} divides by each trial's own healthy rate, we
also divided the fault-window rate of every trial by one reference, the median healthy rate of the
36 trials (0.80 blocks/s). The in-band baseline then has 55--57\% under checkpoint
suspension, 15--18\% under a DA outage and 16--18\% under both; the out-of-band baseline
62\%, 19--20\% and 18--22\%; Engram 96--100\% and the soft baseline 96--100\% under every fault.
The comparison does not depend on the normalization. One further in-band DA trial is not in the table: its
pre-fault window did not run healthy (4 blocks in 60\,s, with proposals already at
round 3), so it measures nothing about the fault, and it was repeated; the
criterion is the pre-fault window alone, and the excluded trial is kept in the
repository. As a sensitivity check, an earlier set of in-band trials, taken before the
Bitcoin receipt tolerance was capped, committed 53--57\% of the healthy rate under checkpoint
suspension and 14--19\% under the DA and combined faults (three trials per cell). The
cap bears on the Bitcoin receipt only, so the DA and combined rates are not expected to depend on it; the
difference to the 17--25\% above is of the size of the variation between trials of one
baseline, and we do not attribute it to the cap.

\emph{After the fault}, the in-band baseline commits nothing in the 120\,s
after the layer returns in all three DA trials, in the combined trials and in
the checkpoint trials. In one earlier trial observed beyond the recovery
window, the chain was still stuck after twelve minutes, with proposals rejected
at round 57 and above. The predicate is satisfiable again once the layer
returns, so the stall comes from where this baseline reads the receipts. The
checkpoint submitter, the DA publisher and the sensor reads run inside
\texttt{ExtendVote}, which ABCI++ calls only when a validator is about to
precommit, and \texttt{ProcessProposal} receives no vote extensions, so the
certificate of a proposal is built from the last committed block. Once
proposals are rejected, no validator precommits and no new receipt is produced,
so the baseline has no path back without an external restart. The out-of-band
baseline reads the layers at each validator whether or not a block commits, and
it resumes in every one of its nine trials: 48--57\,s after the fault ended under
checkpoint suspension, 21--25\,s under a DA outage and 51--74\,s under the
combined fault. We did not separate the two dependencies from each other. The price of the
out-of-band design is that two honest validators can read different values and
disagree on a proposal, which shows in the rounds above and in the longest gap
of about 6.7\,s. All seven validators read the same Bitcoin node and the same
Celestia chain on one host, so their readings are far more alike here than they
would be across independent nodes.
Engram commits 139--151 blocks (94--103\% of its healthy
rate) under every fault, with a longest gap of 1.4\,s under DA faults and about 6\,s
under checkpoint suspension (rounds 2--3), and resumes within two seconds of the fault ending.
The 6\,s gap occurs once in each of the three checkpoint trials, at round 0 and
35--38\,s before the FSM first reaches \textsc{Sovereign}; the soft baseline, which
runs the same FSM, shows it in its three trials as well (Section~\ref{subsec:limitations}).

The withdrawal columns separate the systems on withdrawal exclusion. After the circuit
breaker has locked, Engram committed none; this count is a consistency check of
the breaker, because the window opens when Engram's own FSM reaches
\textsc{Sovereign}. The receipt-gated baselines, which have no breaker, committed 3--7
withdrawals after that moment under DA faults (in-band) or 7--10 (out-of-band) and
none (in-band) or 3 (out-of-band) under checkpoint suspension, few only because they commit few blocks; both
admit withdrawal-marked transactions while degraded, as the soft baseline does, so
the low count after the lock moment reflects how few blocks they commit, not a
safeguard. Engram's cost shows in the
last column: it committed 15--19 withdrawals before locking under DA faults and
55--59 under checkpoint suspension, because it needs about 40\,s and 125\,s to
leave \textsc{Anchored} for \textsc{Sovereign}. These counts follow from the
workload (one withdrawal-marked transaction every 2\,s) and from Bitcoin running
thirty times faster than on mainnet; they are not a production exposure. The
workload-independent measures and a bound are in Section~\ref{sec:discussion}. 

The soft baseline shows the other side of the trade-off. It runs Engram's FSM
and keeps Engram's liveness, 96--100\% of the healthy block rate with a longest
gap of 1.4\,s under DA faults and about 6\,s under checkpoint suspension, but it never locks, so it committed 27--29 withdrawals after
the lock moment under checkpoint suspension and 72--73 under DA faults, against
none for Engram. Counted by the FSM state at each height, 71--97 of its 175
withdrawals under checkpoint suspension and 131--133 of about 180 under DA faults
were committed while the state was \textsc{Sovereign} or \textsc{Recovering}.
Only Engram combines the liveness of the soft chain with the zero count of a
chain that locks.

\textbf{Composite fault.} Engram was then run through a continuous, scripted
composite-fault sequence under a standing wide-area latency profile
(Table~\ref{tab:schedule}): a Bitcoin fault, then a full DA outage (all seven
bridges stopped), then three cycles of packet loss on the validator P2P
network, healed in turn. The
packet loss degrades peer quality without removing quorum connectivity.
Fig.~\ref{fig:fsm-timeline} shows the resulting timeline.

The composite-fault timeline (Fig.~\ref{fig:fsm-timeline}) is in Appendix~\ref{app:extended}.

In the first run, made with the heuristic fallback enabled, over 341 seconds (742 samples, 126 polling ticks), the height
advanced continuously from 905 to 1093 with no stalled height, and all seven
validators agreed on the AppHash at every tick (zero divergent heights, zero
RPC errors). The chain entered \textsc{Sovereign} under the combined fault
and, after the faults were healed, the recovery proof was verified at
$t{=}319$\,s, the chain entered \textsc{Recovering} at $t{=}321$\,s and
\textsc{Anchored} at $t{=}324$\,s. Withdrawals stayed locked throughout
\textsc{Sovereign} and \textsc{Recovering}, as the circuit-breaker invariant
requires. These timings come from a single trial.

We then repeated the sequence five times, each on a freshly deployed chain, with proof-only recovery and an unproven-tail
bound of 128 blocks, so that recovery needed a verified proof. In all five
runs the height advanced throughout (from 373--387 to 747--760), the
validators agreed on the AppHash at every sampled height (zero divergent
heights; one RPC error in one run), and the chain ended in \textsc{Anchored}.
Times are those at which the median (fourth) validator first reported the state, from the first sample. In two of the five runs some validators were briefly in \textsc{Suspicious}, ordinary sensor noise at a threshold of two, before the first fault. The chain reached \textsc{Sovereign} 38--53\,s after the start (mean 44.9\,s, 95\% CI
$\pm$7.2\,s over five runs, Student $t$), four of the five passing through
\textsc{Suspicious} first (16--22\,s). \textsc{Recovering} began 311--324\,s in,
after the faults were healed, and lasted 233--257\,s (mean 240.6\,s, 95\% CI
$\pm$12.0\,s), against 3\,s in the first run; most of that time is the cost of
the recovery proof. In one run two of the seven validators fell back from
\textsc{Recovering} to \textsc{Sovereign} for 2\,s before returning.
We then ran the sequence five times on the cloud testbed (Testbed~B, Appendix~\ref{app:setup}): seven validators in three
regions, proof-only recovery with a tail bound of 128 and two proving workers, each run on a freshly deployed chain, at commit \texttt{3fe656dce021}, in which the sensor reads run in background loops and \texttt{ExtendVote} composes the latest result
(Section~\ref{sec:protocol}), with the times defined as above.

All five runs started with every validator in \textsc{Anchored} and none showed \textsc{Suspicious}. The height advanced throughout (from 218--228 to 461--494), the
validators agreed on the AppHash at every sampled height (zero divergent heights, zero RPC errors), and the chain ended in \textsc{Anchored}.
The median validator left \textsc{Anchored} 33--77\,s after the start (mean 45.2\,s, 95\% CI $\pm$22.6\,s), against a Bitcoin delay injected at 20\,s and the bridges stopped at 94--96\,s; this is the time of the first departure from \textsc{Anchored}
because the passage through \textsc{Sovereign} to \textsc{Recovering} can be shorter than the sampling tick of about 3\,s. In every run six or seven validators fell back from \textsc{Recovering} to \textsc{Sovereign} (the first fall-back of the median validator at 62--104\,s),
between 6 and 13 times in total per run. The Bitcoin delay was removed at 297--315\,s and the bridges were restarted at 301--319\,s. The median validator entered \textsc{Recovering} for the last time at 314--331\,s, and every validator reported \textsc{Anchored} 23--62\,s later (mean 40.5\,s, $\pm$21.8\,s), at 345--385\,s after the start.
The \texttt{ExtendVote} calls of the five runs (25{,}526 calls) had a median of 0.1\,ms in each of the three regions and a maximum of 10.6\,ms.
Each of the five runs logged two proofs rejected at submission by the chain, one within 40\,s of the start and one 368--418\,s after the start, when the chain had already returned to \textsc{Anchored}; recovery completed through the other segments (see Section~\ref{subsec:rq3} for a logged case). All 44 segments the prover logged had both boundary blocks in round 0, with no retry.
Five runs per measure, so the intervals are descriptive.
The seven-scenario sweep (Fig.~\ref{fig:s1-s7}) is in Appendix~\ref{app:extended}.

\subsection{RQ3: Hysteresis and Recovery}
\label{subsec:rq3}

\textbf{Hysteresis and flapping.} Hysteresis should suppress oscillation (``flapping'') without unduly
delaying legitimate transitions. We swept the threshold on each of the
three gated edges -- \textsc{Recovering}$\to$\textsc{Anchored},
\textsc{Anchored}$\to$\textsc{Suspicious}, and
\textsc{Suspicious}$\to$\textsc{Anchored} -- in-process under fixed-seed
synthetic noise and live on the seven-validator cluster under real network
noise (Fig.~\ref{fig:hysteresis}).

The hysteresis figure and table (Fig.~\ref{fig:hysteresis}, Table~\ref{tab:hysteresis-down}) and their discussion are in Appendix~\ref{app:extended}.

\textbf{Recovery: what ends \textsc{Recovering}.} The state exposes whether the proof
condition holds, not which path set it, so we ran one 120\,s DA outage per
setting and read the prover's log. With the heuristic fallback enabled, the
chain left \textsc{Recovering} after 36\,s while the prover's single proof was
rejected as invalid (the interval had grown while it was built): the fallback,
not the proof, ended recovery. With proof-only recovery and the default bound of
8 unproven tail blocks, the chain never left \textsc{Recovering}, not even after
genesis (observed for over thirty minutes in two deployments), although the
chain accepted the prover's proofs (seven in one deployment). A 256-header proof
takes about two minutes to build and submit, more than 90 blocks at 1.2\,s, so
the tail was always longer than the bound. With the bound raised to 128 blocks
the chain left \textsc{Recovering} 341--356\,s after entering it in each of five
runs on fresh chains (mean 347.2\,s, 95\% CI $\pm$8.1\,s), the first valid proof
being accepted 341--353\,s in (recorded in three of the five). The time is
mostly the wait for the prover's batch of 256 headers and its build time, not
the 68--73 blocks of the outage interval. The bound must therefore exceed the proving
latency measured in blocks, and the blocks of that tail are protected by
consensus alone. No block of these intervals was committed in a round other than
0. The RQ2 trials above and the repeated composite runs on Testbed~A use proof-only recovery
with the bound of 128, as do all runs on Testbed~B (the composite, load, hung-RPC, window, delay and boundary-round runs); the first composite run, the seven-scenario runs and the
data-availability runs used the fallback.

\textbf{Recovery with a boundary block committed after round 0.} On Testbed~B (proof-only recovery, tail bound 128, seven validators in three regions) we ran three recoveries from \textsc{Sovereign} on one chain, and repeated the three on a second, freshly deployed chain (Table~\ref{tab:round-recovery}; first run at commit \texttt{fe419f406c46}, second at \texttt{ee271a1c6b2b}, which adds a log line for an accepted proof). Each froze the Bitcoin daemon and stopped the seven Celestia bridges until every validator reported \textsc{Sovereign}, then restored both until every validator reported \textsc{Recovering}. In the second and third recovery two validators adjacent in the proposer rotation were stopped; five of seven keep the quorum, and a height whose first proposer was one of them commits in round 1 or 2. A control file made the prover end a segment only at a block that committed in the wanted round, and hold the finished proof; the prover dropped five headers from the end and retried before it found such a block. While the proof was held we submitted mutated copies of it, then released it.

\begin{table}[tbp]
\caption{Recoveries on Testbed~B with the last boundary block committed in round 0, 1 and 2; each cell gives the first run, then the second (commit rounds read from the chain's \texttt{/commit} for the first and last block of the proof). The recovery with round 0 is the control, with no validator stopped. ``Not informative'': the repeated transaction was byte-identical to the prover's and the mempool answered it.}
\label{tab:round-recovery}
\centering
\footnotesize
\begin{tabular}{@{}lrrr@{}}
\toprule
 & \textbf{Round 0} & \textbf{Round 1} & \textbf{Round 2} \\
\midrule
Heights covered by the held proof & 235--251 / 231--246 & 416--457 / 431--464 & 590--617 / 603--638 \\
Commit round of first / last block & 0 / 0 in both & 0 / 1 in both & 0 / 2 in both \\
Audit-trail pointer on Celestia (height) & none / 62 & 92 / 92 & 121 / 121 \\
Held proof logged as accepted by a validator & not logged / yes & not logged / yes & not logged / yes \\
Seven validators \textsc{Anchored} afterwards & yes / yes & yes / yes & yes / yes \\
Mutated proofs rejected (of 7) & 7 / 7 & 7 / 7 & 7 / 7 \\
Same proof submitted again & not informative / rejected & rejected / rejected & rejected / rejected \\
\bottomrule
\end{tabular}
\end{table}

The chain accepted each held proof and all seven validators returned to \textsc{Anchored}; the chain stayed in \textsc{Recovering} while the proof was held. In the first run the acceptance follows from the prover's count of accepted proofs, from the absence of a failure logged for the segment and from the rejection of the repeated proof; in the second a validator logged the acceptance with the heights of the proof.

Seven mutations of each proof were rejected: a changed $rt_{\text{last}}$, first-block hash, last-block hash, validator-set hash, count or first height, each by its own check before the UltraHonk verification (the validator log names the check), and a changed proof byte by the verification itself. Submitting an accepted proof again was rejected in the chain's handler, at the $rt_{\text{last}}$ check where the validator log named it, in every recovery except the first control, where the repeated transaction was byte-identical to the prover's and the mempool answered it, so that case does not test the chain.

A burst of 24 proofs that pass the binding checks but fail verification, sent in the control recovery, was met by the limit of eight verifications per block: in the first run 13 were refused by the limit and 11 were verified and rejected, in the second 8 and 16 (the number verified depends on how many blocks the burst spans).

A wrong claimed round and a signing power below two thirds cannot be tested on the chain, because no proof exists for them: the circuit does not satisfy its constraints, and the tests of the circuit cover them (Section~\ref{sec:protocol}).

There were two recoveries per round, on two chains. An earlier attempt of the same procedure (two recoveries, rounds 1 and 2, commit \texttt{3fe656dce021}) is not reported: in both, the proof that the prover held was refused when it was released, and in the first the chain had fallen back to \textsc{Sovereign} during the hold. Its harness mutated the newest instead of the oldest held proof, could not tell a repeated transaction from one that reached the handler, and recorded no check names, so the cause of the refusals was not established; the two runs reported here followed it, with the harness and the log corrected.

A segment queued behind the accepted one was rejected after the chain returned to \textsc{Anchored} (in the first recovery, the log of the validators ends with a rejection by the $rt_{\text{new}}$ check), which is consistent with the pruning of the interval's headers on the return and with the two proofs rejected in each composite run above; we did not isolate it.

\textbf{The prover on Testbed~B.} The proving experiment was repeated on the cloud testbed
(Appendix~\ref{app:setup}) with the paper's parameters (Table~\ref{tab:fsm-params}): proof-only recovery and a tail bound of 128. The
prover host handed each proof to one of $N$ worker machines, each building one proof at a time, for $N{=}1,2,3,4$.
A run held the chain in \textsc{Sovereign} for 15 minutes by suspending checkpoint submission and released it; it is one run per
setting. The block rate was \resCloudRate\,blocks/s, and a proof took 76--77\,s from its start to its accepted submission
(12, 22, 23 and 23 proofs accepted). The unproven tail had a 95th percentile of \resCloudTailOne{}, \resCloudTailTwo{}, 47 and 46 blocks
for one to four workers (maxima 63, 49, 52 and 51). The gate was open for 100\% of the samples at the bound of 128 and at 64 in every run;
at 32 it was open for 4--12\% and at 8 for none; the bound of 64 passed with a margin of one block (maximum 63) and is not a safe setting. The percentiles come from 76--77 autocorrelated samples of one hold per setting. The chain returned to \textsc{Anchored} on a verified proof 89--100\,s after the release.
With one and two workers the figures follow $rL(1+1/N)$ (Section~\ref{sec:discussion}), which gives 63 and 47 blocks
for the measured rate and latency; with three and four they stay at the value for two, as Section~\ref{sec:discussion} explains.
They are not comparable in blocks with the Testbed~A runs above, whose block rate was about twice as high.

The data-availability backfill runs and controls (Tables~\ref{tab:backfill} and~\ref{tab:backfill-controls}) and the exploratory adaptive down threshold are in Appendix~\ref{app:extended}.

\subsection{RQ4: Cost}
\label{subsec:rq4}

\begin{table}[tbp]
\caption{Overhead of the extended consensus object against vanilla CometBFT on Testbed~A (seven validators on one host except the one-process column; 20\,tx/s, 60\,s per repeat; mean $\pm$ 95\% CI over the repeats). ``ext.\ on'' is vanilla CometBFT with vote extensions enabled.}
\label{tab:overhead}
\centering
\footnotesize
\begin{tabular}{@{}lrrrr@{}}
\toprule
 & \multicolumn{3}{c}{\textbf{Vanilla}} & \textbf{Engram} \\
\cmidrule(lr){2-4}
\textbf{Metric} & 1 process & 7 val., no ext. & 7 val., ext.\ on & 7 val. \\
\midrule
Repeats (runs) & 5 & 5 & 3 & 5 \\
Rounds per height & 1.000 & 1.000 & 1.000 & 1.000 \\
Proposals rejected / height & 0 & 0 & 0 & 0 \\
Nil precommits / height & 0 & 0 & 0 & 0 \\
Block interval, mean (s) & 1.034 $\pm$ 0.001 & \resBlockVanilla{} $\pm$ 0.015 & 1.229 $\pm$ 0.023 & \resBlockEngram{} $\pm$ 0.020 \\
Block interval p50 / p95 (s) & 1.032 / 1.042 & 1.220 / 1.313 & 1.217 / 1.300 & 1.241 / 1.335 \\
Blocks/s & 0.967 & 0.811 & 0.814 & 0.793 \\
Committed tx/s & 20.00 & 20.00 & 20.00 & 20.00 \\
Tx commit latency p50 / p95 (s) & 0.67 / 1.16 & 0.99 / 1.57 & 0.99 / 1.56 & 1.03 / 1.63 \\
Mempool backlog, mean / max & 8.2 / 24 & 14.5 / 31 & 14.6 / 31 & 15.5 / 44 \\
Tx per block (mean) & 17.6 & 21.1 & 21.1 & 21.5 \\
\bottomrule
\end{tabular}
\\[2pt]
\begin{minipage}{\linewidth}
\raggedright
\footnotesize Per-height values come from scanning \emph{every} committed
height (commit round, nil precommits, block time). Nil prevotes are not stored
on chain; rejected proposals (distinct height/round/reason in the validators'
logs) are the proxy for why a validator would prevote nil. ``ext.\ on'' keeps
vote extensions enabled in the genesis with empty extensions; ``no ext.'' is a
fresh genesis without them. The Engram arm is a fresh deployment; a first
deployment gave 1.242\,$\pm$\,0.005\,s (three runs). Transaction counts exclude
the health certificate in \texttt{Txs[0]}. Latency runs from submission to the
first 0.25\,s poll that sees the transaction in a block. The 1-process arm is a
local single node.
\end{minipage}
\end{table}

The certificate-scaling measurements (Table~\ref{tab:cert-scaling}) and the analysis of where the cost is and is not are in Appendix~\ref{app:extended}.

The load sweep up to 500\,tx/s on one host (Table~\ref{tab:throughput}) is in Appendix~\ref{app:extended}; on one host Engram's latency grows with load where vanilla's does not.

\begin{table}[tbp]
\caption{Load of 100\,tx/s for 60\,s on Testbed~B (mean $\pm$ 95\% Student-$t$ interval over three freshly deployed chains per column, one measurement each). Vanilla ran at commit \texttt{6e010cd} on 12 October; the other three columns at commit \texttt{3fe656d} on 16--17 October. ``Empty ext.'' signs a constant healthy observation without reading any sensor; ``sync'' reads the Bitcoin and Celestia nodes inside \texttt{ExtendVote}; ``snapshot'' composes the latest background reading. The share in \textsc{Anchored} is the fraction of the samples of one validator's FSM state during the measurement.}
\label{tab:testbedb-load}
\centering
\footnotesize
\begin{tabular}{@{}lrrrr@{}}
\toprule
 & \textbf{Vanilla} & \textbf{Empty ext.} & \textbf{Sync} & \textbf{Snapshot} \\
\midrule
Block interval (s) & 1.433 $\pm$ 0.027 & 1.447 $\pm$ 0.023 & 2.773 $\pm$ 0.209 & 1.453 $\pm$ 0.030 \\
Latency p50 (s) & 2.70 $\pm$ 0.03 & 3.09 $\pm$ 1.51 & 3.74 $\pm$ 0.32 & 2.70 $\pm$ 0.11 \\
Latency p95 (s) & 3.61 $\pm$ 0.02 & 5.82 $\pm$ 9.30 & 5.49 $\pm$ 0.80 & 3.63 $\pm$ 0.10 \\
Latency p99 (s) & 3.85 $\pm$ 0.04 & 6.14 $\pm$ 9.71 & 6.44 $\pm$ 1.15 & 3.88 $\pm$ 0.16 \\
\texttt{ExtendVote}, mean (ms) & n/a & 0.01 $\pm$ 0.00 & 924.25 $\pm$ 30.82 & 0.12 $\pm$ 0.01 \\
Rounds per height & 1.000 $\pm$ 0.000 & 1.000 $\pm$ 0.000 & 1.000 $\pm$ 0.000 & 1.000 $\pm$ 0.000 \\
Share of samples in \textsc{Anchored} & n/a & 1.000 $\pm$ 0.000 & 0.231 $\pm$ 0.597 & 0.923 $\pm$ 0.191 \\
Busiest validator CPU (\%) & 23.8 $\pm$ 5.7 & 33.8 $\pm$ 34.1 & 20.5 $\pm$ 17.6 & 34.9 $\pm$ 32.5 \\
Committed / offered & 1.000 $\pm$ 0.000 & 1.000 $\pm$ 0.000 & 1.000 $\pm$ 0.000 & 1.000 $\pm$ 0.000 \\
\bottomrule
\end{tabular}

\end{table}

\textbf{Load on separate machines.} On Testbed~B, where each validator has its own machine, we offered 100\,tx/s for 60\,s to four configurations on three freshly deployed chains each (Table~\ref{tab:testbedb-load}), with the load generator on the operator's machine, outside the cloud: vanilla CometBFT, which does not run the vote-extension code; vanilla's code with an empty vote extension on every vote and a constant healthy observation (empty); Engram with the sensor reads inside \texttt{ExtendVote} (sync); and Engram with background sensor reads (snapshot). The client's round trip to the validators and a queueing term that grows with the block interval enter the latency, so only the block interval is independent of the client. Every configuration committed every offered transaction in every deployment, with one round per height.

The block interval was $1.433 \pm 0.027$\,s for vanilla, $1.447 \pm 0.023$\,s for empty, $1.453 \pm 0.030$\,s for snapshot and $2.773 \pm 0.209$\,s for sync. The mean \texttt{ExtendVote} time was $924 \pm 31$\,ms for sync and $0.12 \pm 0.01$\,ms for snapshot ($0.01$\,ms for empty). The 95th-percentile latency was $3.61 \pm 0.02$\,s for vanilla, $3.63 \pm 0.10$\,s for snapshot and $5.49 \pm 0.80$\,s for sync. For empty it was 3.65 and 3.67\,s in two deployments and 10.14\,s in the first, whose block interval (1.446\,s), rounds per height and committed rate were those of the others; we did not find the cause of that one deployment, and its busiest validator used 50\% of its CPU against 26\% in the other two. With three deployments per column the comparison is descriptive.

The FSM did not stay in \textsc{Anchored} in every deployment. With sync the share of samples in \textsc{Anchored} was 0.50, 0.15 and 0.04 (\textsc{Suspicious} in all three, and \textsc{Recovering} and \textsc{Sovereign} in the third); with snapshot it was 0.85, 1.00 and 0.92 (\textsc{Suspicious} in two); with empty it was 1.00 by construction. The latency of the sync column is therefore the latency of a chain that spent most of the measurement outside \textsc{Anchored}, and a run at this load does not say how long a chain stays in \textsc{Anchored}. The busiest validator used 13--50\% of its CPU in all four configurations, with no relation to the configuration.

Two readings follow. First, the empty configuration has the block interval of vanilla, so the extra vote-extension round does not account for the interval of sync on this network. Second, the interval of snapshot is that of vanilla as well, which is consistent with the sensor reads inside \texttt{ExtendVote} being the cause: a precommit quorum of five of seven cannot form without a validator of ap-southeast-1, and the \texttt{ExtendVote} of sync has to cross to the Bitcoin and Celestia nodes in us-west-2. The runs do not report \texttt{ExtendVote} per region and do not isolate that effect.

\textbf{Slow and hung Bitcoin RPC.} We froze the Bitcoin daemon with \texttt{docker pause} for 240\,s on a chain in \textsc{Anchored}, so that every validator's connection opened and every request waited. Table~\ref{tab:testbedb-hang} gives the two builds. With inline reads, \texttt{ExtendVote} took $1.61$\,s on average during the freeze (maximum 1.78\,s), twice the 800\,ms bound of one RPC call, which fits a read that chains two calls, and the block interval rose from 2.39 to 2.98\,s. With background reads it stayed at $0.06$\,ms (maximum 0.6\,ms) and the interval was 1.42\,s, that of the healthy phase. Both chains kept committing, left \textsc{Anchored} for \textsc{Sovereign} and returned to \textsc{Anchored} 30.8\,s (inline) and 27.6\,s (background) after the daemon resumed. The background build reported a Bitcoin gap of 8 (the worst case, the sovereign threshold) in all 1{,}106 snapshot lines after the freeze took hold, and none with a gap of 0: a frozen layer was not reported as healthy. The freshness flag was never raised in this test, because each call returns within its 800\,ms bound and the background reader therefore always delivered a fresh reading of \emph{unreachable}; the branch for a reading older than the window is covered by a unit test and by the window sweep below. One run per build.

\begin{table}[tbp]
\caption{Hung Bitcoin RPC on Testbed~B (daemon frozen for 240\,s, one run per build, no transaction load). Block intervals from the sampled heights; states are shares of the samples of all seven validators.}
\label{tab:testbedb-hang}
\centering
\scriptsize
\resizebox{\linewidth}{!}{\begin{tabular}{@{}llrrrl@{}}
\toprule
\textbf{Build} & \textbf{Phase} & \textbf{Blocks} & \textbf{Interval (s)} & \texttt{ExtendVote} mean / p95 / max (ms) & \textbf{States (share of samples)} \\
\midrule
Inline reads & 60\,s healthy & 23 & 2.39 & 896.66 / 1221.9 / 2285.6 & Anchored 100\% \\
 & 240\,s hung & 80 & 2.98 & 1613.85 / 1719.3 / 1779.2 & Sovereign 99\%, Anchored 1\% \\
 & until \textsc{Anchored} & 12 & 2.13 & 977.63 / 1613.1 / 2290.6 & Recovering 45\%, Anchored 29\%, Sovereign 26\% \\
\midrule
Background reads & 60\,s healthy & 44 & 1.36 & 0.07 / 0.1 / 5.0 & Anchored 84\%, Suspicious 16\% \\
 & 240\,s hung & 167 & 1.42 & 0.06 / 0.1 / 0.6 & Sovereign 99\%, Anchored 1\% \\
 & until \textsc{Anchored} & 19 & 1.21 & 0.06 / 0.1 / 0.3 & Recovering 60\%, Anchored 36\%, Sovereign 5\% \\
\bottomrule
\end{tabular}
}
\end{table}

\textbf{\texttt{ExtendVote} against injected delay.} To vary the latency of the Bitcoin RPC instead of freezing it, \texttt{tc netem} added a fixed delay to every response that left the Bitcoin host, for 90\,s per value (0, 100, 300 and 500\,ms; Table~\ref{tab:testbedb-latency}), on a freshly deployed chain of each build and with no transaction load. With inline reads the mean \texttt{ExtendVote} time rose from 0.87\,s to 1.64, 3.09 and 4.65\,s, about 7.5\,ms for each millisecond of delay, and the block interval from 2.54 to 3.32, 4.43 and 5.82\,s. With background reads the mean was 0.06--0.07\,ms at every value (maxima 0.3--1.9\,ms) and the block interval 1.36--1.65\,s. At 500\,ms the added delay and the round trip to the Bitcoin host exceeded the 800\,ms bound of one call, and the background build reported the layer as failed: its validators left \textsc{Anchored} for \textsc{Sovereign} and \textsc{Recovering} (79\% of the samples in \textsc{Recovering}), while the validators with inline reads were in \textsc{Suspicious} for 61\% of the samples. One 90\,s window per value and build, with 11--50 blocks per window; the intervals are descriptive, and samples that failed are not counted in the states.

\begin{table}[tbp]
\caption{\texttt{ExtendVote} time and block interval against the delay added to every Bitcoin RPC response (Testbed~B, 90\,s per value, one chain per build, no transaction load). The interval is the span of the highest sampled height over the window after the first 15\,s.}
\label{tab:testbedb-latency}
\centering
\scriptsize
\resizebox{\linewidth}{!}{\begin{tabular}{@{}llrrrl@{}}
\toprule
\textbf{Build} & \textbf{Added delay (ms)} & \textbf{Blocks} & \textbf{Interval (s)} & \texttt{ExtendVote} mean / p95 / max (ms) & \textbf{States (share of samples)} \\
\midrule
Inline reads & 0 & 28 & 2.54 & 865.89 / 1200.1 / 2230.2 & Anchored 73\%, Suspicious 27\% \\
 & 100 & 22 & 3.32 & 1644.33 / 2833.4 / 3538.8 & Anchored 100\% \\
 & 300 & 14 & 4.43 & 3087.39 / 4584.3 / 6134.2 & Anchored 100\% \\
 & 500 & 11 & 5.82 & 4646.79 / 8507.8 / 8924.8 & Suspicious 61\%, Anchored 39\% \\
\midrule
Background reads & 0 & 47 & 1.45 & 0.07 / 0.1 / 0.3 & Anchored 86\%, Suspicious 14\% \\
 & 100 & 50 & 1.36 & 0.06 / 0.1 / 0.3 & Anchored 100\% \\
 & 300 & 40 & 1.65 & 0.06 / 0.1 / 1.9 & Anchored 100\% \\
 & 500 & 50 & 1.42 & 0.06 / 0.1 / 0.4 & Recovering 79\%, Anchored 15\%, Sovereign 6\% \\
\bottomrule
\end{tabular}
}
\end{table}

\textbf{Freshness window.} We repeated the hung-RPC test with the window of the snapshot rule set to 1, 2, 5 and 10\,s (\texttt{ENGRAM\_SENSOR\_MAX\_AGE\_MS}; Table~\ref{tab:testbedb-maxage}), each on a freshly deployed chain, with 60\,s healthy and 240\,s frozen. In none of the four runs did a validator report a Bitcoin gap of 0 while the daemon was frozen (0 of 4{,}508 snapshot lines), and \texttt{ExtendVote} took 0.06\,ms on average in each. The flag that marks a stale or missing reading was raised in the frozen phase only with the window of 1\,s (9 of 1{,}148 lines) and in the healthy phase only with that window (4 of 301 lines); with 2, 5 and 10\,s it was never raised. The reader loop runs every 0.5\,s and one call may wait 0.8\,s, so a window of 1\,s can be shorter than the age of the latest reading, and from 2\,s on it was not in these runs. The time to return to \textsc{Anchored} after the daemon resumed was 42.1, 29.0, 38.8 and 27.6\,s, with no order in the window. In the healthy phase the FSM was in \textsc{Suspicious} for 21, 17, 19 and 0\% of the samples, which the flag does not explain (it was raised in none of the lines of the 2 and 5\,s runs); we did not isolate the cause. One run per window. The run with 5\,s repeats the test of the previous paragraphs with the same settings (27.6\,s; flag never raised).

\begin{table}[tbp]
\caption{Hung Bitcoin RPC with the freshness window of the snapshot rule set to 1, 2, 5 and 10\,s (Testbed~B, one freshly deployed chain per window, daemon frozen for 240\,s). ``Degraded'' counts the \texttt{sensor\_snapshot} lines that carry the flag for a missing or stale reading; ``healthy reads'' counts the lines of the frozen phase that report a Bitcoin gap of 0.}
\label{tab:testbedb-maxage}
\centering
\scriptsize
\resizebox{\linewidth}{!}{\begin{tabular}{@{}rrrrrr@{}}
\toprule
\textbf{Window (s)} & \textbf{Degraded, healthy} & \textbf{Degraded, frozen} & \textbf{Healthy reads, frozen} & \texttt{ExtendVote} mean / max (ms) & \textbf{Back to \textsc{Anchored} (s)} \\
\midrule
1 & 4 / 301 & 9 / 1148 & 0 & 0.06 / 0.4 & 42.1 \\
2 & 0 / 308 & 0 / 1141 & 0 & 0.06 / 0.7 & 29.0 \\
5 & 0 / 301 & 0 / 1127 & 0 & 0.06 / 2.0 & 38.8 \\
10 & 0 / 294 & 0 / 1092 & 0 & 0.06 / 0.4 & 27.6 \\
\bottomrule
\end{tabular}
}
\end{table}
